% GENERATED FILE. Edit canonical lessons, then run npm run generate:books.
% canonical-source-hash: fnv1a64:46ee941b1aa73320
% canonical-lessons: FR-C257-geographie, FR-C257-mathematiques, FR-C257-art, FR-C257-equipe, FR-C257-anniversaire, FR-R257-first-pass-words-from-chapters-251-253, FR-R257-second-pass-words-from-chapters-254-257

\chapter{Geography, Mathematics, Art, Team, Birthday}
\label{ch:fr-geography-mathematics-art-team-birthday}
\hlchaptermodality{257}

\begin{chapteropening}
\textbf{By the end of this chapter:} \emph{I can say geography, mathematics, art, a team and a birthday.}

Complete the last lesson of chapter 257: I can say geography, mathematics, art, a team and a birthday.
\end{chapteropening}

\section[la géographie]{la géographie — geography}

\label{lesson:FR-C257-geographie}



\begin{quote}
Before the new one: say the French for a violin, then the French for a drum.
\end{quote}

\subsection*{You'll want to know: la géographie}
\textbf{la géographie} — \textquotedblleft{}geography\textquotedblright{}.

\begin{grammarlens}[title={What you've built}]
One more word for this chapter.
\end{grammarlens}

\subsection*{Your turn}
\begin{itemize}
\raggedright
  \item \emph{Say it:} \emph{la géographie}
  \item \emph{Say it:} \emph{la géographie}, once more
  \item \emph{Say it:} say \emph{la géographie}
  \item \emph{Recall:} say the French for a violin, then the French for a drum, then say \emph{la géographie} again
\end{itemize}

\begin{tcolorbox}[breakable,colback=teal!4,colframe=teal!35!black,title={Before you move on}]
What does la géographie mean? (\textquotedblleft{}Geography\textquotedblright{}.) Say it once more.
\end{tcolorbox}
\section[les mathématiques]{les mathématiques — mathematics}

\label{lesson:FR-C257-mathematiques}



\begin{quote}
Before the new one: say the French for a drum, then the French for geography.
\end{quote}

\subsection*{You'll want to know: les mathématiques}
\textbf{les mathématiques} — \textquotedblleft{}mathematics\textquotedblright{}.

\begin{grammarlens}[title={What you've built}]
One more word for this chapter.
\end{grammarlens}

\subsection*{Your turn}
\begin{itemize}
\raggedright
  \item \emph{Say it:} \emph{les mathématiques}
  \item \emph{Say it:} \emph{les mathématiques}, once more
  \item \emph{Say it:} say \emph{les mathématiques}
  \item \emph{Recall:} say the French for a drum, then the French for geography, then say \emph{les mathématiques} again
\end{itemize}

\begin{tcolorbox}[breakable,colback=teal!4,colframe=teal!35!black,title={Before you move on}]
What does les mathématiques mean? (\textquotedblleft{}Mathematics\textquotedblright{}.) Say it once more.
\end{tcolorbox}
\section[l'art]{l'art — art}

\label{lesson:FR-C257-art}



\begin{quote}
Before the new one: say the French for geography, then the French for mathematics.
\end{quote}

\subsection*{You'll want to know: l'art}
\textbf{l'art} — \textquotedblleft{}art\textquotedblright{}.

\begin{grammarlens}[title={What you've built}]
One more word for this chapter.
\end{grammarlens}

\subsection*{Your turn}
\begin{itemize}
\raggedright
  \item \emph{Say it:} \emph{l'art}
  \item \emph{Say it:} \emph{l'art}, once more
  \item \emph{Say it:} say \emph{l'art}
  \item \emph{Recall:} say the French for geography, then the French for mathematics, then say \emph{l'art} again
\end{itemize}

\begin{tcolorbox}[breakable,colback=teal!4,colframe=teal!35!black,title={Before you move on}]
What does l'art mean? (\textquotedblleft{}Art\textquotedblright{}.) Say it once more.
\end{tcolorbox}
\section[l'équipe]{l'équipe — a team}

\label{lesson:FR-C257-equipe}



\begin{quote}
Before the new one: say the French for mathematics, then the French for art.
\end{quote}

\subsection*{You'll want to know: l'équipe}
\textbf{l'équipe} — \textquotedblleft{}a team\textquotedblright{}.

\begin{grammarlens}[title={What you've built}]
One more word for this chapter.
\end{grammarlens}

\subsection*{Your turn}
\begin{itemize}
\raggedright
  \item \emph{Say it:} \emph{l'équipe}
  \item \emph{Say it:} \emph{l'équipe}, once more
  \item \emph{Say it:} say \emph{l'équipe}
  \item \emph{Recall:} say the French for mathematics, then the French for art, then say \emph{l'équipe} again
\end{itemize}

\begin{tcolorbox}[breakable,colback=teal!4,colframe=teal!35!black,title={Before you move on}]
What does l'équipe mean? (\textquotedblleft{}A team\textquotedblright{}.) Say it once more.
\end{tcolorbox}
\section[l'anniversaire]{l'anniversaire — a birthday}

\label{lesson:FR-C257-anniversaire}



\begin{quote}
Before the new one: say the French for art, then the French for a team.
\end{quote}

\subsection*{You'll want to know: l'anniversaire}
\textbf{l'anniversaire} — \textquotedblleft{}a birthday\textquotedblright{}.

\begin{grammarlens}[title={What you've built}]
One more word for this chapter.
\end{grammarlens}

\subsection*{Your turn}
\begin{itemize}
\raggedright
  \item \emph{Say it:} \emph{l'anniversaire}
  \item \emph{Say it:} \emph{l'anniversaire}, once more
  \item \emph{Say it:} say \emph{l'anniversaire}
  \item \emph{Recall:} say the French for art, then the French for a team, then say \emph{l'anniversaire} again
\end{itemize}

\begin{tcolorbox}[breakable,colback=teal!4,colframe=teal!35!black,title={Before you move on}]
What does l'anniversaire mean? (\textquotedblleft{}A birthday\textquotedblright{}.) Say it once more.
\end{tcolorbox}
\section[(dialogue)]{Words from chapters 251-253 — a first pass}

\label{lesson:FR-R257-first-pass-words-from-chapters-251-253}



\begin{quote}
Say the words for a team and a birthday.
\end{quote}

\subsection*{Your turn}
Say these aloud:

\begin{itemize}
\raggedright
  \item a result, a difference, an example, the truth, a lie
  \item a tax, a form, an envelope, a stamp, a parcel
  \item a driving licence, a helmet, a tunnel, a building site, a block of flats
\end{itemize}

\begin{tcolorbox}[breakable,colback=teal!4,colframe=teal!35!black,title={Before you move on}]
A result? (\textbf{le résultat}.) A difference? (\textbf{la différence}.) An example? (\textbf{l'exemple}.) The truth? (\textbf{la vérité}.) A lie? (\textbf{le mensonge}.) A tax? (\textbf{l'impôt}.) A form? (\textbf{le formulaire}.) An envelope? (\textbf{l'enveloppe}.) A stamp? (\textbf{le timbre}.) A parcel? (\textbf{le colis}.) A driving licence? (\textbf{le permis de conduire}.) A helmet? (\textbf{le casque}.) A tunnel? (\textbf{le tunnel}.) A building site? (\textbf{le chantier}.) A block of flats? (\textbf{l'immeuble}.)
\end{tcolorbox}
\section[(dialogue)]{Words from chapters 254-257 — a second pass}

\label{lesson:FR-R257-second-pass-words-from-chapters-254-257}



\begin{quote}
Say the words for a trunk and a thorn.
\end{quote}

\subsection*{Your turn}
Say these aloud:

\begin{itemize}
\raggedright
  \item a trunk, a thorn, a lawn, a path, a cave
  \item rubbish, recycling, energy, electricity, gas
  \item a tool, a guitar, a piano, a violin, a drum
  \item geography, mathematics, art, a team, a birthday
\end{itemize}

\begin{tcolorbox}[breakable,colback=teal!4,colframe=teal!35!black,title={Before you move on}]
A trunk? (\textbf{le tronc}.) A thorn? (\textbf{l'épine}.) A lawn? (\textbf{la pelouse}.) A path? (\textbf{le sentier}.) A cave? (\textbf{la grotte}.) Rubbish? (\textbf{les déchets}.) Recycling? (\textbf{le recyclage}.) Energy? (\textbf{l'énergie}.) Electricity? (\textbf{l'électricité}.) Gas? (\textbf{le gaz}.) A tool? (\textbf{l'outil}.) A guitar? (\textbf{la guitare}.) A piano? (\textbf{le piano}.) A violin? (\textbf{le violon}.) A drum? (\textbf{le tambour}.) Geography? (\textbf{la géographie}.) Mathematics? (\textbf{les mathématiques}.) Art? (\textbf{l'art}.) A team? (\textbf{l'équipe}.) A birthday? (\textbf{l'anniversaire}.)
\end{tcolorbox}
